\paragraph{\trVanishingTheorem。}\label{sub:vanishing-theorem}
上文提到的最后一个\trCompositionLaw{}在我们接下来的讨论中还要用到多次，因此值得把它表述得稍更精确一些。为简单起见，设 $\pi_k\colon E_k\to S_k$（$k=1,2$）是两个\trLefschetzFibration，第一个带\trStriplikeEnd，第二个的基\trCompact，二者各配有一个（可能移动的）\trLagrange\trBoundaryCondition{}$F_k$。此外，我们还要有\trPointedBoundary{}$z_1\in\partial S_1$、$z_2\in\partial S_2$，使得两根\trFiber{}$(E_k)_{z_k}$ 都能等同于同一个 $M$，而诸 $(F_k)_{z_k}$ 则成为同一个\trSubmanifold{}$L\subset M$。我们还假定，在这两个\trMarkedPoint{}附近，两个\trFibration{}都是局部平凡的，而两个\trBoundaryCondition{}都是\trLocallyConstant{}的。选取\trRelative\trPerturbationDatum{}$(K_k,J_k)$，它们在 $(E_k)_{z_k}$ 附近是平凡的（$K_k=0$，$J_k=I_{E_k}$）。现在在我们的\trMarkedPoint{}处作出 $S_1$ 与 $S_2$ 的\trBoundaryConnectedSum。这一步可以用两种方式来想：或者先把诸 $z_k$ 去掉，再在所得的\trPunctured\trSurface{}上引入\trStriplikeEnd，然后照 \eqref{eq:striplike-gluing} 那样对这些端执行一个标准的\trGluing{}过程，取某个\trGluingLength{}$l$；或者考察以 $z_k$ 为中心、以半\trDisk{}$D_+=D\cap H$ 为模型的局部坐标，然后使用形状为 $z\mapsto\rho z^{-1}$ 的过渡映射，其中 $\rho\in(-1;0)$。这两幅图像是完全\trEquivalence{}的（其中有通常的关系式 $l=-\log(\rho)/\pi$）。我们之所以提到这两种图像，是因为它们各有其直观之处，并且因为它们暗示了两个略有不同、但同样可行的分析框架：第一个导向通常\trFloer{}理论的一个 \trMorseBott{}版本，而第二个则接近（相对）\trGromov--\trWitten{}理论中\trHolomorphicDisc{}的\trGluing{}程序。无论如何，所得结果上带有一个\trLefschetzFibration{}$\pi\colon E\to S$，连同随之而来的\trBoundaryCondition{}$F$ 与\trRelative\trPerturbationDatum{}$(K,J)$；而基本的\trGluing{}定理 [131, Equation (2.28)] 断言：若适当的\trRegularity{}条件得到满足，并且 $l$ 充分大（或 $\rho$ 充分接近 $0$），则
\begin{equation}\label{eq:fibre-product-gluing}
\scr{M}_{E/S}(\{y_\zeta\})^0\cong
\coprod_{p+q=n}\scr{M}_{E_1/S_1}(\{y_\zeta\})^p\times_L\scr{M}_{E_2/S_2}^q.
\end{equation}
\begingroup\tolerance=9999\emergencystretch=3em\hbadness=10000
这里，上标表示维数：在右边，它们经过了调整，使得关于\trEvaluationMap{}$\mathrm{ev}_{z_1}\colon\scr{M}_{E_1/S_1}(\{y_\zeta\})\to(F_1)_{z_1}\cong L$ 与 $\mathrm{ev}_{z_2}\colon\scr{M}_{E_2/S_2}\to(F_2)_{z_2}\cong L$ 所取的\trFibreProduct{}是零维的（上文提到的\trRegularity{}假设包含了 $\mathrm{ev}_{z_1}$、$\mathrm{ev}_{z_2}$ 相互处于\trGeneralPosition{}这条性质）。现在我们把它转述成关于相应\trInvariant{}的一个命题。对 $E_2$，考察由 $\mathrm{ev}_{z_2}$ 所代表的\trHomology{}类，并把它记作 $\Phi_{E_2/S_2,z_2}\in H_*(L;\KK)$；对 $E$ 则直接取 \eqref{eq:cphi-fibration-map}。
\begin{lemma}\label{thm:fibre-product-vanishing}
若 $\Phi_{E_2/S_2,z_2}$ 为零，则 $\Phi_{E/S}$ 亦为零。
\end{lemma}
\endgroup
其证明是直截了当的：取一条以 $\mathrm{ev}_{z_2}$ 为边缘的（伪）\trCycle{}$b\colon B\to L$。此时，$\scr{M}_{E_1/S_1}(\{y_\zeta\})\times_L B$ 的零维部分给出\trFloerCochainGroup{}之间的一个映射，它就是 $\Phi_{E/S}$ 的一个零伦。
\begin{remark}\label{thm:cochain-level-vanishing}
在我们对这个引理的全部应用中，其中的\trCompact{}部分都将取如\cref{thm:phi-es-vanishing}中那样。此时还可以附加 $2c_1(M,L)=0$ 这一假设，从而得到一个稍更精细的结果。当 $n>1$ 时，空间 $\scr{M}_{E_2/S_2}$ 的维数是 $2n-1>n$，因此\trFibreProduct{} \eqref{eq:fibre-product-gluing} 为空。在剩下的情形 $n=1$ 中，唯一的非平凡项是 $p=0$、$q=1$ 的那一项。由\cref{thm:phi-es-vanishing}我们已经知道\trEvaluationMap
\begin{equation*}
\mathrm{ev}_{z_2}\colon\scr{M}_{E_2/S_2}^1\longrightarrow L
\end{equation*}
是偶\trDegree{}的。于是，对 $\scr{M}_{E_1/S_1}(\{y_\zeta\})^0$ 中的每个\trIsolated{}点 $u_1$，$\mathrm{ev}_{z_1}(u_1)$ 都是 $\mathrm{ev}_{z_2}$ 的\trRegularValue（由\trTransversality{}假设），从而满足 $\mathrm{ev}_{z_2}(u_2)=\mathrm{ev}_{z_1}(u_1)$ 的映射 $u_2$ 有偶数个。这表明该\trFibreProduct{}同样含有偶数个点。结论是：对这个特定选取的\trPerturbationDatum，上链层面的\trInvariant{}$C\Phi_{E/S}$ 本身即为零。
\end{remark}
